\chapter{Herpesvirus Infections}
\index{Herpesvirus Infections}

\section*{Definition and Overview}
Herpesvirus infections are caused by the eight human herpesviruses: herpes simplex virus types 1 and 2 (HSV-1, HSV-2), varicella-zoster virus (VZV), Epstein--Barr virus (EBV), cytomegalovirus (CMV), human herpesvirus 6 (HHV-6), human herpesvirus 7 (HHV-7), and Kaposi sarcoma-associated herpesvirus (HHV-8). These viruses share the ability to remain latent in the body after the first infection and reactivate later. They cause a wide range of illnesses, from common cold sores and chickenpox to glandular fever, serious infections in immunocompromised people, and some cancers.

\section*{Epidemiology}
Herpesviruses are among the most common human infections. Over 90\% of adults have antibodies to HSV-1, EBV, CMV, and HHV-6. Chickenpox affected almost every local child before vaccination, and shingles affects one in three adults. EBV causes most cases of glandular fever in young people. CMV is the most common congenital infection. In people with HIV or organ transplants, herpesviruses cause significant disease. Vaccines are available for chickenpox and shingles, and antivirals treat several herpesvirus infections.

\section*{Aetiology and Pathophysiology}
\begin{itemize}
    \item \textbf{HSV-1 and HSV-2:} cold sores and genital herpes, transmitted by contact, latent in nerve ganglia
    \item \textbf{Varicella-zoster virus:} chickenpox (primary) and shingles (reactivation), latent in nerve roots
    \item \textbf{Epstein--Barr virus:} glandular fever, transmitted by saliva, latent in B lymphocytes; linked to some cancers
    \item \textbf{Cytomegalovirus:} usually silent in healthy people; causes serious disease in the fetus and immunocompromised people, latent in white cells
    \item \textbf{HHV-6 and HHV-7:} cause roseola in infants, latent for life
    \item \textbf{HHV-8:} causes Kaposi sarcoma, mainly in immunocompromised people
    \item \textbf{Common themes:} primary infection, lifelong latency, reactivation with immune decline, and transmission by body fluids
\end{itemize}

\section*{Clinical Features}
\begin{itemize}
    \item Cold sores and genital sores from HSV
    \item Chickenpox: widespread itchy blisters in children
    \item Shingles: painful banded rash in older adults
    \item Glandular fever: fever, sore throat, swollen glands, and fatigue from EBV
    \item CMV mononucleosis: similar to EBV but less throat involvement
    \item Roseola: fever followed by a rash in infants
    \item In immunocompromised people: pneumonia, hepatitis, colitis, retinitis, and encephalitis
    \item Congenital CMV: hearing loss, developmental delay, and birth defects
\end{itemize}

\section*{Differential Diagnosis}
\begin{itemize}
    \item Other viral infections cause fever, rash, sore throat, and fatigue
    \item Bacterial tonsillitis mimics EBV and must be excluded before antibiotics
    \item HIV and other immunodeficiencies are diagnosed in at-risk patients
    \item Specific tests identify each herpesvirus
    \item Serology, PCR, and antigen tests distinguish primary infection from reactivation
\end{itemize}

\section*{When to Seek Medical Care}
See a doctor for persistent fever, severe sore throat, unusual rashes, or symptoms in pregnancy or with a weakened immune system. Urgent care is needed for neurological symptoms, eye involvement, or severe illness in vulnerable people.

\textbf{Contact your local emergency services for:}
\begin{itemize}
    \item Severe headache, confusion, or seizures (possible encephalitis)
    \item Eye pain or vision changes with a herpes infection
    \item A newborn with fever, rash, or lethargy
    \item Severe breathing difficulty or chest pain
    \item High fever with a very unwell child or immunocompromised adult
\end{itemize}

\section*{Diagnosis and Investigations}
\begin{itemize}
    \item \textbf{Viral PCR:} the most sensitive test, on swabs, blood, urine, or cerebrospinal fluid
    \item \textbf{Serology:} antibodies distinguish past infection from current infection
    \item \textbf{Viral culture and antigen tests:} older methods still used in some settings
    \item \textbf{Full blood count:} atypical lymphocytes in EBV and CMV
    \item \textbf{Monospot test:} rapid screening for EBV in glandular fever
    \item \textbf{Imaging and CSF analysis:} for neurological complications
    \item \textbf{Newborn testing:} for congenital CMV in affected babies
\end{itemize}

\section*{Management}
\begin{itemize}
    \item \textbf{Supportive care:} most herpesvirus infections in healthy people resolve without specific treatment
    \item \textbf{Antivirals:} aciclovir and valaciclovir for HSV and VZV; ganciclovir and valganciclovir for CMV
    \item \textbf{Pain and fever relief:} paracetamol and other supportive measures
    \item \textbf{Avoid antibiotics:} they do not treat viruses and can cause rash in EBV
    \item \textbf{Immunocompromised care:} prophylaxis, early treatment, and dose adjustment
    \item \textbf{Pregnancy management:} monitoring and treatment of congenital CMV and neonatal herpes
    \item \textbf{Vaccination:} chickenpox and shingles vaccines prevent VZV disease
    \item \textbf{Cancer treatment:} specific therapy for Kaposi sarcoma and EBV-associated cancers
\end{itemize}

\section*{Complications}
\begin{itemize}
    \item Post-herpetic neuralgia after shingles
    \item Herpes encephalitis and keratitis
    \item Glandular fever complications: splenic rupture (rare), liver inflammation, and prolonged fatigue
    \item Congenital CMV: hearing loss and developmental disability
    \item Serious disseminated disease in immunocompromised people
    \item EBV- and HHV-8-associated cancers in at-risk groups
    \item Transmission to contacts and newborns
\end{itemize}

\section*{Prognosis and Outcomes}
Most herpesvirus infections in healthy people run a self-limiting course and resolve without treatment, though the viruses remain in the body for life. Primary EBV and CMV infections cause significant illness but recover fully in nearly all cases. The main risks occur in newborns and immunocompromised people, where antiviral therapy and preventive care have greatly improved outcomes. Vaccination has virtually eliminated chickenpox complications and markedly reduces shingles. For most people, herpesvirus infections are a manageable part of life.

\section*{Prevention and Self-Care}
\begin{itemize}
    \item Vaccination: chickenpox and shingles vaccines are safe and effective
    \item Practise hand hygiene and avoid sharing drinks and utensils
    \item Cover herpes sores and shingles rashes to prevent spread
    \item Avoid contact between pregnant women and people with active infections where relevant
    \item Rest, fluids, and simple analgesia for viral illness
    \item Immunocompromised people should follow specialist infection-prevention advice
    \item Seek early medical care for symptoms in pregnancy, in infants, or with immune weakness
    \item Reduce stress and maintain general health to limit reactivation
\end{itemize}

\section*{Sources and Further Reading}
The following reputable sources provide further information for healthcare students and patients seeking reliable, current advice about herpesvirus infections.

\begin{itemize}
    \item Healthdirect. \textit{Herpesviruses and related infections.} https://www.healthdirect.gov.au/
    \item Australian Government Department of Health. \textit{Chickenpox and shingles immunisation.} https://www.health.gov.au/
    \item CMV Australia and CMV education for women. \textit{Congenital CMV information.} https://www.cmvaustralia.org.au/
    \item NHS. \textit{Epstein--Barr virus and glandular fever.} https://www.nhs.uk/conditions/glandular-fever/
    \item Mayo Clinic. \textit{Infection information for cytomegalovirus.} https://www.mayoclinic.org/
    \item Centers for Disease Control and Prevention. \textit{Herpesvirus disease information.} https://www.cdc.gov/
\end{itemize}
